\DSource{0135}{46}
\hypertarget{AB01-PDF0135-P01}{}
\DLine{AB01-PDF0135-body-L001}{}{}{gnoscitur reliquum latus AM trianguli, quod est [cosinus] arcus FH et anguli FAH (1). Igi-}
\DLine{AB01-PDF0135-body-L002}{}{}{tur est etiam linea EM nota. Cum triangulum FME rectangulum sit, eius hypotenusa EF}
\DLine{AB01-PDF0135-body-L003}{}{}{cognoscitur, ex qua linea FM deprehenditur; arcus autem super FM est arcus anomaliae.}
\hypertarget{AB01-PDF0135-P02}{}
\DLine{AB01-PDF0135-body-L004}{}{}{\Indent Si, ut posuimus, arcus FH $30^\circ$ complectitur, eius sinus erit 30 partium, e partibus qua-}
\DLine{AB01-PDF0135-body-L005}{5}{}{rum semidiametrus AF 60 continet; sed ea ratione qua linea AF $2^p4'45''$ complectitur, erit}
\DLine{AB01-PDF0135-body-L006}{}{}{linea FM $1^p2'22''\,\qfrac{1}{2}$ (2), reliqua igitur AM $1^p48'2''$, et linea EM (3) $61^p48'2''$ (4); unde}
\DLine{AB01-PDF0135-body-L007}{}{}{patet [lineam EF] $61^p48'35''$ circiter esse ({\itshape d}). Sed ex partibus quarum linea EF 60 tantum}
\DLine{AB01-PDF0135-body-L008}{}{}{habet, linea FM $1^p0'33''$ (5) continebit, et arcus super eam positus $0^\circ57'49''$ circiter; et}
\DLine{AB01-PDF0135-body-L009}{}{}{haec est quantitas anomaliae motus Solis, scilicet arcus HF. Est igitur arcus TA eclipticae}
\DLine{AB01-PDF0135-body-L010}{10}{}{$29^\circ2'11''$, qui antea $30^\circ$ erat; nam epicycli centrum a T ad A motum est ut Sol in epicyclo}
\DLine{AB01-PDF0135-body-L011}{}{}{a H ad F.}
\hypertarget{AB01-PDF0135-P03}{}
\DLine{AB01-PDF0135-body-L012}{}{}{\Indent Nunc epicyclum GQP consideremus, cuius centrum est B; locum Solis G ponamus, sit-}
\DLine{AB01-PDF0135-body-L013}{}{p. 70.}{que arcus QG (6), quem Sol a puncto Q apogei percurrit, $150^\circ$ (7); remanent igitur $30^\circ$ arcui}
\DLine{AB01-PDF0135-body-L014}{}{}{PG, qui a loco Solis ad punctum perigei porrigitur. Lineam EG et perpendicularem GK du-}
\DLine{AB01-PDF0135-body-L015}{15}{}{camus; patet triangula BKG et GKE rectangula esse, et latera BG, BE haud ignota; nam}
\DLine{AB01-PDF0135-body-L016}{}{}{BG est semidiametrus epicycli (8), et BE semidiametrus circuli eclipticae. Angulus et arcus}
\DLine{AB01-PDF0135-body-L017}{}{}{GP dati sunt; perpendicularis GK et reliqua linea BK notae; [ergo etiam lineae KE et EG}
\par\noindent\begin{minipage}[t]{0.49\FullSourceWidth}\vspace{0pt}\centering\hypertarget{AB01-PDF0135-F01}{}\includegraphics[width=\linewidth]{AB01-PDF0135-F01.png}\end{minipage}\hfill\begin{minipage}[t]{0.465\FullSourceWidth}\vspace{0pt}\setlength{\GutterOffset}{0.5349999999999999\FullSourceWidth}
\DLine{AB01-PDF0135-body-L018}{}{}{cognoscuntur] (9). Cum arcus GP, ut posui-}
\DLine{AB01-PDF0135-body-L019}{}{}{mus, $30^\circ$ contineat, eius sinus quoque $30^p$}
\DLine{AB01-PDF0135-body-L020}{20}{}{erit; arcus autem complementaris, qui su-}
\DLine{AB01-PDF0135-body-L021}{}{}{per KB est, sinum habebit $51^p57'41''$ (10).}
\DLine{AB01-PDF0135-body-L022}{}{}{--- At ex partibus quarum linea BG $2^p4'\,\qfrac{3}{4}$}
\DLine{AB01-PDF0135-body-L023}{}{}{habet, perpendicularis KG $1^p2'22''\,\qfrac{1}{2}$ (11)}
\DLine{AB01-PDF0135-body-L024}{}{}{continebit; remanet linea KB $1^p48'2''$, qua}
\DLine{AB01-PDF0135-body-L025}{25}{}{re EK fere $58^p11'58''$, et EG fere $58^p12'$}
\DLine{AB01-PDF0135-body-L026}{}{}{$34''$ (12) continebit ({\itshape e}). Sed ex partibus qua-}
\DLine{AB01-PDF0135-body-L027}{}{}{rum EG 60 habet, cathetus KG $1^p4'17''$ (13)}
\DLine{AB01-PDF0135-body-L028}{}{}{complectetur, et arcus super eam positus}
\DLine{AB01-PDF0135-body-L029}{}{}{$1^\circ1'24''$ (14), iuxta rationem qua circulus}
\DLine{AB01-PDF0135-body-L030}{30}{}{triangulum rectangulum BKG circumdans}
\DLine{AB01-PDF0135-body-L031}{}{}{in $360^\circ$ dividitur. Et hic est arcus anoma-}
\DLine{AB01-PDF0135-body-L032}{}{}{liae, scilicet arcus PG. Qua re arcus NB (15)}
\DLine{AB01-PDF0135-body-L033}{}{}{eclipticae, quem cognoscere volebamus, $31^\circ$}\end{minipage}\par
\DLine{AB01-PDF0135-body-L034}{}{}{$1'24''$ (16) complectetur.}
\par\vspace{6pt}\hrule height.25pt\vspace{5pt}
\noindent\begin{minipage}[t]{.48\textwidth}\vspace{0pt}\fontsize{9}{11.5}\selectfont\setlength{\GutterOffset}{0pt}
\hypertarget{AB01-PDF0135-N01}{}
\DLine{AB01-PDF0135-notes_left-L001}{}{}{\Indent (1) Cod.: « quod est quod angulo FAH et arcui}
\DLine{AB01-PDF0135-notes_left-L002}{}{}{« FH deest ad quadrantem perficiendum ».}
\hypertarget{AB01-PDF0135-N02}{}
\DLine{AB01-PDF0135-notes_left-L003}{}{}{\Indent (2) Pro $22''$ Plato habet $55''$. Persaepe in Pla-}
\DLine{AB01-PDF0135-notes_left-L004}{}{}{tonis editione pro 2 legitur 5, qui error codicibus}
\DLine{AB01-PDF0135-notes_left-L005}{}{}{Latinis vel editoris incuriae tribuendus est.}
\hypertarget{AB01-PDF0135-N03}{}
\DLine{AB01-PDF0135-notes_left-L006}{}{}{\Indent (3) Codex EF.}
\hypertarget{AB01-PDF0135-N04}{}
\DLine{AB01-PDF0135-notes_left-L007}{}{}{\Indent (4) Plato verba seqq. omittit et pro $2''$ habet $35''$.}
\hypertarget{AB01-PDF0135-N05}{}
\DLine{AB01-PDF0135-notes_left-L008}{}{}{\Indent (5) Plato $1^\circ33'$.}
\hypertarget{AB01-PDF0135-N06}{}
\DLine{AB01-PDF0135-notes_left-L009}{}{}{\Indent (6) Cod. QM.}
\hypertarget{AB01-PDF0135-N07}{}
\DLine{AB01-PDF0135-notes_left-L010}{}{}{\Indent (7) Male Plato 120.}
\hypertarget{AB01-PDF0135-N08}{}
\DLine{AB01-PDF0135-notes_left-L011}{}{}{\Indent (8) Cod.: eclipticae.}
\hypertarget{AB01-PDF0135-N09}{}
\DLine{AB01-PDF0135-notes_left-L012}{}{}{\Indent (9) Addidi, Platonem sequens; sed additio non}
\DLine{AB01-PDF0135-notes_left-L013}{}{}{est necessaria.}
\end{minipage}
\hfill\begin{minipage}[t]{.48\textwidth}\vspace{0pt}\fontsize{9}{11.5}\selectfont\setlength{\GutterOffset}{0pt}
\hypertarget{AB01-PDF0135-N10}{}
\DLine{AB01-PDF0135-notes_right-L001}{}{}{\Indent (10) Partes apud Plat. $59^p$.}
\hypertarget{AB01-PDF0135-N11}{}
\DLine{AB01-PDF0135-notes_right-L002}{}{}{\Indent (11) Deest $\qfrac{1}{2}$ in cod.; Plato pro $22''$ habet $55''$}
\DLine{AB01-PDF0135-notes_right-L003}{}{}{(cfr. adnot. 2).}
\hypertarget{AB01-PDF0135-N12}{}
\DLine{AB01-PDF0135-notes_right-L004}{}{}{\Indent (12) Plato $58^p7'34''$.}
\hypertarget{AB01-PDF0135-N13}{}
\DLine{AB01-PDF0135-notes_right-L005}{}{}{\Indent (13) Secundae apud Platonem $13''$.}
\hypertarget{AB01-PDF0135-N14}{}
\DLine{AB01-PDF0135-notes_right-L006}{}{}{\Indent (14) Plato $1^p4'54''$.}
\hypertarget{AB01-PDF0135-N15}{}
\DLine{AB01-PDF0135-notes_right-L007}{}{}{\Indent (15) Cod. IB (\textarabic{ى} I pro \textarabic{ن} N); Plato GB. In figura}
\DLine{AB01-PDF0135-notes_right-L008}{}{}{codicis et Platonis punctum N lineaque NT desunt;}
\DLine{AB01-PDF0135-notes_right-L009}{}{}{punctum T vero, ut supra notavi, perperam prope}
\DLine{AB01-PDF0135-notes_right-L010}{}{}{locum intersectionis epicycli HF et eclipticae collo-}
\DLine{AB01-PDF0135-notes_right-L011}{}{}{catum.}
\hypertarget{AB01-PDF0135-N16}{}
\DLine{AB01-PDF0135-notes_right-L012}{}{}{\Indent (16) Minuta in cod. $4'$; Plat. $1^\circ4'54''$.}
\end{minipage}
